\documentclass[
  aps,            % American Physical Society style
  prresearch,     % journal substyle (Phys. Rev. Research); change to prb/prx/prl as desired
  reprint,        % two-column published format; use "preprint" for the one-column manuscript style
  superscriptaddress,
  floatfix,
  nofootinbib
]{revtex4-2}

% ---------------------------------------------------------------------------
%  Hybrid resolution refinement + rodeo algorithm on the pairing Hamiltonian
%  Built from the Jupyter notebook ResolutionRefPairing.ipynb, the EIGEN-Q
%  research proposal, and Qian et al., Eur. Phys. J. A 60, 151 (2024).
%  RevTeX 4-2 / APS style.
% ---------------------------------------------------------------------------

\usepackage{amsmath,amssymb}
\usepackage{physics}
\usepackage{graphicx}
\usepackage{hyperref}

\newcommand{\Hlow}{H_{\mathrm{low}}}
\newcommand{\Hhigh}{H_{\mathrm{high}}}
\newcommand{\Eobj}{E_{\mathrm{obj}}}

\begin{document}

\title{Hybrid resolution refinement and the rodeo algorithm for
eigenstate preparation: a demonstration on the pairing Hamiltonian}

\author{Scott Bogner}
\affiliation{Facility for Rare Isotope Beams and Department of Physics and Astronomy, Michigan State University, East Lansing, MI 48824, USA}
\author{Cecilie Glittum}
\affiliation{Department of Physics, University of Oslo, N-0316 Oslo, Norway}
\author{Heiko Hergert}
\affiliation{Facility for Rare Isotope Beams and Department of Physics and Astronomy, Michigan State University, East Lansing, MI 48824, USA}
\author{Morten Hjorth-Jensen}
\email{mhjensen@uio.no}
\affiliation{Department of Physics, University of Oslo, N-0316 Oslo, Norway}
\author{Gunnar Lange}
\affiliation{Department of Physics, University of Oslo, N-0316 Oslo, Norway}
\author{Ryan LaRose}
\affiliation{Department of Computational Mathematics, Science and Engineering and Facility for Rare Isotope Beams, Michigan State University, East Lansing, MI 48824, USA}
\author{Dean Lee}
\affiliation{Facility for Rare Isotope Beams and Department of Physics and Astronomy, Michigan State University, East Lansing, MI 48824, USA}
\author{Francesco Pietro Massel}
\affiliation{Department of Science and Industry Systems, University of South-Eastern Norway, N-3179 Borre, Norway}
\affiliation{Department of Physics, University of Oslo, N-0316 Oslo, Norway}

\date{\today}

\begin{abstract}
The efficient preparation of correlated eigenstates is a central bottleneck for
quantum simulation: variational quantum eigensolvers (VQE) suffer from
exponentially vanishing gradients, adiabatic state preparation requires evolution
times that diverge when the spectral gap closes, and projective filters such as
quantum phase estimation and the rodeo algorithm converge exponentially only when
the input state already has a sizeable overlap with the target eigenstate.  We
demonstrate, end to end and on a single controlled testbed --- the constant-pairing
Hamiltonian --- a hybrid pipeline that resolves these tensions: a cheap
low-resolution VQE state is prolonged into a larger model space, adiabatically
\emph{refined} across resolutions, and finally polished by the rodeo algorithm.
Resolution refinement keeps the path gapped, so the adiabatic time scales only as
$T\sim 1/\Delta E$ and the overlap $p$ with the high-resolution ground state rises
toward unity; the rodeo algorithm then removes the residual error exponentially,
with a post-selected infidelity $1-\mathcal{F}_M\lesssim \tfrac{1-p}{p}\,2^{-M}$ in
the number of cycles $M$ and a post-selection acceptance that stays of order $p$.
For four particles refined up to four levels ($k=4$, eight qubits) the refined input
($p\simeq0.998$) is filtered to a fidelity $1-\mathcal{F}\sim10^{-6}$ within one or
two cycles, while an unrefined prolonged guess ($p\simeq0.87$) still reaches
$\sim10^{-4}$ in a handful of cycles.  All gate-level building blocks (UCCSD state
preparation, the prolongation embedding, and the ancilla-based rodeo cycle) are
implemented explicitly and validated against exact diagonalisation.  A
gate-count analysis quantifies the depth advantage of the hybrid: a single
UCCSD--VQE preparation at $k=4$ already costs $1312$ CNOTs and must be
re-executed throughout a classical optimisation loop, whereas one rodeo cycle
costs a few hundred CNOTs per Trotter step, carries no variational parameters,
and converges exponentially.  The
construction is the pairing-model realisation of the EIGEN-Q proposal and a concrete
template for hardware demonstrations on mid-circuit-measurement platforms such as the
Magne neutral-atom quantum computer.
\end{abstract}

\maketitle



% ===========================================================================
\section{Introduction}
% ===========================================================================
Determining the low-lying eigenvalues and eigenvectors of a many-body Hamiltonian is
one of the grand challenges of quantum many-body theory, where the Hilbert-space
dimension grows exponentially with the number of constituents.  Quantum computers
offer new routes to eigenvalue estimation and eigenstate preparation, but the
state-of-the-art methods each face an obstacle at scale.  The variational quantum
eigensolver (VQE)~\cite{Peruzzo2014} is limited by barren plateaus --- gradients that
vanish exponentially with system size~\cite{McClean2018}.  Adiabatic state
preparation requires evolution times that diverge if the gap between the tracked
state and the rest of the spectrum closes along the path~\cite{Albash2018}.
Projective spectral filters --- quantum phase estimation~\cite{Abrams1999} and the
rodeo algorithm (RA)~\cite{Choi2021,Qian2024} --- converge exponentially in the
number of cycles, but only once the initial state has a significant squared overlap
$p$ with the target eigenstate; their cost grows as $1/p$, and $p$ itself typically
falls exponentially with system size.

These limitations are complementary, which suggests a hybrid strategy.
\emph{Resolution refinement}~\cite{Bogner2026} bootstraps an eigenstate across
model-space scales: an eigenstate of a coarse (low-resolution) Hamiltonian is lifted
by a prolongation operator into a finer space and adiabatically evolved to the
fine-resolution eigenstate.  Because refinement preserves the low-energy spectral
structure, the path stays gapped and the required adiabatic time scales only as the
inverse physical gap (times a slowly growing system-size factor), delivering an
output with $p\approx 1$.  That output is exactly the high-overlap input the rodeo
algorithm needs to converge exponentially.  This is the core idea of the EIGEN-Q
pipeline: (i) prepare a low-resolution eigenstate; (ii) prolong and adiabatically
refine to a final time $T^\ast$; (iii) apply the rodeo algorithm to reach the target
fidelity and observables with exponential convergence.

In this work we realise the full pipeline end to end on a transparent and exactly
solvable testbed, the constant-pairing Hamiltonian familiar from nuclear and
condensed-matter many-body theory~\cite{LNP936}.  Working on a model whose exact
spectrum is available at every stage lets us validate each component quantitatively:
the gate-level UCCSD coarse solver against full configuration interaction (FCI) and
coupled-cluster doubles (CCD); the prolongation and adiabatic refinement against the
exact high-resolution eigenstate; and the rodeo filter against the analytic
success-probability formulas of Ref.~\cite{Qian2024}.  All circuits are implemented
explicitly as rotation and CNOT gates (and an ancilla-based rodeo cycle), so the
demonstration doubles as a reference implementation for hardware.

% ===========================================================================
\section{The pairing model and its qubit encoding}
% ===========================================================================
We study the constant-pairing Hamiltonian
\begin{widetext}
\begin{equation}
  H=\delta\sum_{p=0}^{k-1} p\sum_{\sigma=\uparrow\downarrow} n_{p\sigma}
   -\frac{g}{2}\sum_{p,q=0}^{k-1} a^\dagger_{p\uparrow}a^\dagger_{p\downarrow}
     a_{q\downarrow}a_{q\uparrow},
  \label{eq:H}
\end{equation}
\end{widetext}
with $k$ doubly degenerate single-particle levels of spacing $\delta$ (set to
$\delta=1$) and a pairing strength $g$.  The interaction scatters time-reversed
pairs between levels; it conserves the seniority (the number of unpaired particles),
so the ground state of $N$ paired particles lives entirely in the seniority-zero
subspace.  We map level $p$ and spin $\sigma$ to qubit $2p+\sigma$ (so a $k$-level
problem uses $2k$ qubits) and encode the fermionic operators through the
Jordan--Wigner transformation.  The Hartree--Fock (HF) reference fills the lowest
$N/2$ levels, with $E_{\rm HF}(N{=}4,g)=2-g$.

Throughout we use $N=4$ particles as the running example.  The fixed-$N$ sector
dimension is $\binom{2k}{N}$, growing from $1$ at $k=2$ (a single filled
determinant) to $70$ at $k=4$ and $1820$ at $k=8$; the seniority-zero ground-state
problem is correspondingly compact, which makes the model an ideal multi-resolution
laboratory.  ``Basis refinement'' here means \emph{increasing $k$ at fixed $N$}: the
number of single-particle levels (and qubits) grows while the particle number is
held fixed.

% ===========================================================================
\section{The hybrid pipeline}
% ===========================================================================

\subsection{Coarse-state preparation: gate-level UCCSD--VQE}
\label{sec:uccsd}
The coarse eigenstate is prepared with a unitary coupled-cluster ansatz restricted
to singles and doubles (UCCSD), Trotterised into a product of excitation circuits
acting on the HF reference,
\begin{equation}
  \ket{\Psi(\boldsymbol\theta)}
  =\prod_{\mu} e^{\theta_\mu(\hat\tau_\mu-\hat\tau_\mu^\dagger)}\ket{\rm HF}.
\end{equation}
Through the Jordan--Wigner map each anti-Hermitian generator becomes a sum of
commuting Pauli strings, so every factor compiles into a basis change, a CNOT ladder,
a single $R_z(\theta)$ rotation, and the uncompute; a single excitation yields two
Pauli-rotation sub-circuits and a double excitation yields eight.  We optimise with
the exact analytic gradient of the product ansatz (the naive two-term parameter shift
is incorrect beyond the first factor, because later factors acquire frequency-one
cross terms).  Because the pairing force does not couple to seniority-changing
single excitations, all singles amplitudes vanish to machine precision and the
solver is effectively a \emph{doubles} theory, mirroring the classical CCD structure.

\subsection{The prolongation operator}
The prolongation $P$ embeds a low-resolution state into the high-resolution space.
For basis refinement it maps every many-body basis state of the $k_{\rm low}$ space to
the \emph{same} occupation pattern in the $k_{\rm high}$ space, leaving the new
single-particle levels empty.  In canonical Jordan--Wigner ordering this is
gate-free; for lattice refinement the analogous operator distributes each occupied
coarse site over the neighbouring fine sites with two single-qubit rotations and two
CNOTs per site per spatial direction, i.e.\ $O(N_{\rm sites})$ gates per
dimension~\cite{Bogner2026}.  The prolonged state $\ket{\Phi_0}=P\ket{\psi_{\rm low}}$
satisfies $\expval{\Hhigh}{\Phi_0}=E_{\rm low}$ exactly, and already provides a
high-quality first guess for the high-resolution ground state.

\subsection{Resolution refinement (adiabatic phase)}
\label{sec:refine}
The prolonged state is refined by adiabatically following the interpolated,
energy-shifted Hamiltonian
\begin{widetext}
\begin{equation}
  H(s)=\cos^2\!\Big(\tfrac{\pi s}{2}\Big)\,P(\Hlow-\mu)P^\dagger
      +\sin^2\!\Big(\tfrac{\pi s}{2}\Big)\,(\Hhigh-\mu),
  \qquad s=\frac{t}{T},
  \label{eq:path}
\end{equation}
\end{widetext}
evolved by a Suzuki--Trotter product over total time $T\le T^\ast$.  The shift $\mu$
(here $\mu=E_{\rm low}+0.6$) places the tracked low-energy states below the large
null space of $P^\dagger$, so the level followed from $s=0$ to $s=1$ is the
high-resolution ground state.  The instantaneous spectrum stays smooth and gapped
along the entire path; consequently the adiabatic time scales as $T\sim 1/\Delta E$
rather than the $O(1/\Delta E_{\min}^2)$ of generic adiabatic preparation, and the
overlap $p=\abs{\braket{\Psi_{\rm high}}{\Phi(T)}}^2$ rises toward unity.

\subsection{The rodeo algorithm}
\label{sec:rodeo}
The rodeo algorithm~\cite{Choi2021,Qian2024} removes the residual error left by the
adiabatic tail.  One cycle acts on the system register plus a single ancilla
(reused across cycles by mid-circuit measurement): a Hadamard on the ancilla, the
controlled evolution $e^{-iH t_k}$ of the system, a phase gate $P(E\,t_k)$ on the
ancilla (which effectively shifts $H\to H-E$), a second Hadamard, and an ancilla
measurement.  Post-selecting the ancilla in $\ket{0}$ applies the non-unitary filter
$\tfrac12\big(\mathbb{1}+e^{iE t_k}e^{-iH t_k}\big)$ to the system: an eigenstate
$\ket{E_j}$ acquires the factor $e^{i(E-E_j)t_k/2}\cos[(E_j-E)t_k/2]$.  For an
eigenstate input the probability that all $M$ ancilla measurements return $0$ is
therefore~\cite{Qian2024}
\begin{equation}
  P_{0^{M}}\!\big(E\mid\{t_k\}\big)=\prod_{k=1}^{M}
  \cos^2\!\Big[\tfrac{t_k}{2}\big(\Eobj-E\big)\Big],
  \label{eq:rodeo1}
\end{equation}
and averaging the $t_k$ over a Gaussian of RMS width $\sigma$ gives
\begin{equation}
  P_{0^{M}}(E)=\Bigg[\frac{1+e^{-(\Eobj-E)^2\sigma^2/2}}{2}\Bigg]^{M}.
  \label{eq:rodeo2}
\end{equation}
On resonance ($E=\Eobj$) the success probability is unity at every cycle, while
off-resonant weight is suppressed geometrically toward the background $2^{-M}$ (the
arithmetic mean of $\cos^2$ is $1/2$).  Scanning $E$ thus locates eigenvalues as
peaks whose width $\sim 1/\sigma$ sets the energy resolution; fixing $E$ at the
target eigenvalue turns the algorithm into a state-preparation filter.

The controlled evolution is the only system-size-dependent ingredient.  On hardware
$e^{-iH t}$ is realised by a Suzuki--Trotter decomposition over the $2k$
qubits~\cite{Trotter1959,Suzuki1976,Childs2021}; for the single-qubit benchmark of
Ref.~\cite{Qian2024} it reduces to a $Z$--$Y$--$Z$ Euler rotation
$R_{\hat n}(\theta)=e^{-i\theta\,\hat n\cdot\vec\sigma/2}$ implemented through the
Qiskit native gate.  In the simulations below we apply the controlled evolution
exactly within the conserved $N$-particle sector, isolating the algorithmic behaviour
from Trotter error (quantified in Sec.~\ref{sec:discussion}).

\paragraph{Convergence criterion.}
With $E$ tuned to the target eigenvalue $E_0$, the all-zero acceptance probability of
an input $\ket{\psi}=\sum_k c_k\ket{E_k}$ after $M$ cycles is
\begin{equation}
  P_M=p+\sum_{k\neq0}\abs{c_k}^2
  \Big[\tfrac12\big(1+e^{-(E_k-E_0)^2\sigma^2/2}\big)\Big]^{M},
  \quad p=\abs{c_0}^2 .
  \label{eq:PM}
\end{equation}
The target weight $p$ is preserved at every round while off-resonant components decay
geometrically, so $P_M\to p$ and the post-selected fidelity obeys
\begin{equation}
  \mathcal{F}_M=\frac{p}{P_M}\simeq\Big[1+\tfrac{1-p}{p}\,2^{-M}\Big]^{-1},
  \qquad
  1-\mathcal{F}_M\lesssim\frac{1-p}{p}\,2^{-M}.
  \label{eq:fid}
\end{equation}
Reaching fidelity $1-\delta$ therefore needs only
$M\gtrsim\log_2[(1-p)/(p\,\delta)]$ cycles, and since the expected number of
repetitions to obtain an accepted run is $\sim 1/P_M\to 1/p$, the total
controlled-evolution effort scales as $O\!\big(\abs{\log\delta}/(p\,\epsilon)\big)$
at energy resolution $\epsilon$~\cite{Choi2021}.  This is efficient precisely when
resolution refinement supplies $p\approx 1$.

% ===========================================================================
\section{Results}
% ===========================================================================

\subsection{Coarse solver: UCCSD versus exact references}
\label{sec:res-uccsd}
Table~\ref{tab:uccsd} lists the ground-state energies at $g=1$ for $N=4$ as the basis
grows from $k=2$ to $k=8$.  The gate-level UCCSD circuit reproduces the operator-level
(statevector) UCCSD to machine precision and stays variational ($E_{\rm UCCSD}\ge
E_{\rm FCI}$), tracking FCI to $\sim10^{-3}$--$10^{-2}$, whereas CCD over-binds.
Figure~\ref{fig:methods} compares all solvers (exact, CCD, statevector and gate-level
UCCSD, and a hardware-efficient ansatz~\cite{Kandala2017}) as a function of $k$.  The hardware-efficient
ansatz is exact at $k=2$ and reaches FCI at $k=3$ only after a multi-start
optimisation, and struggles at $k=4$ --- the symmetry-blind landscape that motivates
the physically structured UCCSD circuit.  Figure~\ref{fig:scaling} shows that the same
UCCSD solver remains a tight variational bound all the way to $k=8$ (sixteen qubits),
at the cost of a CNOT count that grows steeply ($1312$ CNOTs already at $k=4$); the
compiled circuits and their depths are analysed in Sec.~\ref{sec:circuits}.

\begin{table}[t]
\caption{\label{tab:uccsd}Ground-state energies at $g=1$, $N=4$, for the
constant-pairing Hamiltonian.  FCI is the exact $N$-sector result; CCD is
coupled-cluster doubles; UCCSD is the gate-level circuit of Sec.~\ref{sec:uccsd}.
``\#CNOT'' is the two-qubit gate count of the UCCSD state preparation.}
\begin{ruledtabular}
\begin{tabular}{cccccc}
$k$ & qubits & FCI & CCD & UCCSD & \#CNOT\\
\colrule
2 & 4  & $1.000000$ & $1.000000$ & $1.000000$ & 0\\
3 & 6  & $0.794697$ & $0.794697$ & $0.794697$ & 272\\
4 & 8  & $0.635548$ & $0.630443$ & $0.636987$ & 1312\\
6 & 12 & $0.372664$ & $0.350058$ & $0.378719$ & --\\
8 & 16 & $0.145559$ & $0.097263$ & $0.157856$ & --\\
\end{tabular}
\end{ruledtabular}
\end{table}

\begin{figure*}[t]
\centering
\includegraphics[width=\linewidth]{figs/fig4.png}
\caption{Ground-state energy at $N=4$, $g=1$ obtained with all solvers (left) and
their deviation from FCI (right), versus the number of levels $k$.  The structured
UCCSD circuits track FCI from above; CCD over-binds; the hardware-efficient ansatz
matches FCI for $k\le3$ but loses accuracy at $k=4$.}
\label{fig:methods}
\end{figure*}

\begin{figure*}[t]
\centering
\includegraphics[width=\linewidth]{figs/fig3.png}
\caption{Scalability of the Trotterised UCCSD--VQE: ground-state energy (left) and
deviation from FCI (right) from $k=2$ to $k=8$ ($N=4$, $g=1$).  The same solver
remains variational and tight as the basis --- and hence the qubit count --- grows.}
\label{fig:scaling}
\end{figure*}

\subsection{Resolution refinement: overlap and energy}
\label{sec:res-refine}
At $N=4$ the low-resolution ($k=2$) space is a single filled determinant, so the
prolonged state is the high-resolution HF reference, with first-guess overlaps
$0.917$ for refinement to $k=3$ and $0.868$ to $k=4$.  Adiabatic evolution along the
gapped path of Eq.~\eqref{eq:path} drives both quantities toward their exact targets:
the overlap rises to $\simeq0.999$ and the energy falls from $E_{\rm HF}=1$ to the FCI
value from above --- $0.7964$ for $k=3$ (FCI $0.794697$) and $0.6397$ for $k=4$ (FCI
$0.635548$) at $T=30$ (Fig.~\ref{fig:fourparticle}).  The instantaneous spectrum stays
open along the whole path (Fig.~\ref{fig:refine}, right), confirming the favourable
$T\sim 1/\Delta E$ timescale; the larger $k=4$ space starts lower in overlap and
converges more slowly, exactly as expected from the system-size dependence of the
refinement.

\begin{figure*}[t]
\centering
\includegraphics[width=\linewidth]{figs/fig5.png}
\caption{Resolution refinement $k=2\to k=4$.  Left: overlap of the evolved state with
the exact high-resolution ground state versus the total adiabatic time $T$, for $N=2$
and $N=4$.  Right: the instantaneous spectrum along the interpolation parameter
$s=t/T$ stays smooth and gapped, the condition for efficient refinement.}
\label{fig:refine}
\end{figure*}

\begin{figure*}[t]
\centering
\includegraphics[width=\linewidth]{figs/fig6.png}
\caption{Four-particle refinement up to $k=4$: overlap (left) and physical energy
$\expval{\Hhigh}{\Phi(T)}$ (right) versus $T$ for $k_{\rm high}=3,4$.  The energy
approaches the FCI value (dotted) from above while the overlap approaches unity, both
on a $\sim1/\Delta E$ timescale.}
\label{fig:fourparticle}
\end{figure*}

\subsection{Rodeo filtering of the refined state}
\label{sec:res-rodeo}
We verified the rodeo primitive in two ways: the explicit ancilla circuit agrees with
the compact filter $\tfrac12(\mathbb{1}+e^{iEt}e^{-iHt})$ to $10^{-16}$, and the
all-zero probabilities reproduce Eqs.~\eqref{eq:rodeo1}--\eqref{eq:rodeo2}, with the
Gaussian-averaged single-cycle factor matching $\tfrac12(1+e^{-(\Delta E\,\sigma)^2/2})$
to three digits.

The state fed to the filter is the resolution-refined $k=4$ state of
Sec.~\ref{sec:res-refine}, with overlap $p=0.998$ with the exact ground state; for
contrast we also use the unrefined prolonged guess $\Phi_0$ ($p=0.868$).  Scanning the
target energy $E$ (Fig.~\ref{fig:scan}) produces a peak at the ground-state energy; a
coarse scan ($\sigma=2$) resolves it from the excited band, and a fine scan
($\sigma=8$) locates it at $E=0.6352$ against the exact $E_0=0.635548$.  Starting from
the refined state, essentially all spectral weight sits in the single ground-state
peak; starting from $\Phi_0$, the small residual weight on excited states is visible as
secondary structure above the $2^{-M}$ background.

Fixing $E=E_0$ and applying $M$ cycles (Fig.~\ref{fig:convergence}), the post-selected
fidelity rises exponentially and the cumulative acceptance settles at $\approx p$, in
quantitative agreement with Eqs.~\eqref{eq:PM}--\eqref{eq:fid}.  The refined input
reaches an infidelity $1-\mathcal{F}\simeq5\times10^{-6}$ within one to two cycles;
even the cruder prolonged input reaches $\simeq1.3\times10^{-4}$ in a few cycles, while
its acceptance holds at $0.868$.  The dotted reference lines $\tfrac{1-p}{p}2^{-M}$
bracket the measured infidelities, confirming the convergence criterion.  This is the
concrete pay-off of the pipeline: resolution refinement supplies the high overlap that
turns the rodeo algorithm from a $1/p$-expensive filter into a few-cycle polish.

\begin{figure*}[t]
\centering
\includegraphics[width=\linewidth]{figs/fig7.png}
\caption{Rodeo energy scan ($M=4$).  Left: coarse scan ($\sigma=2$); the dotted
vertical lines mark exact eigenvalues and the dashed line the $2^{-M}$ background.
Right: fine scan ($\sigma=8$) on the ground-state peak, located at $E=0.6352$ (exact
$E_0=0.6355$).}
\label{fig:scan}
\end{figure*}

\begin{figure*}[t]
\centering
\includegraphics[width=\linewidth]{figs/fig8.png}
\caption{Rodeo eigenstate preparation at $E=E_0$, averaged over random time sets.
Left: post-selected infidelity $1-\mathcal{F}_M$ versus the number of cycles $M$
(dotted: the bound $\tfrac{1-p}{p}2^{-M}$).  Right: the post-selection acceptance
converges to the input overlap $p$.}
\label{fig:convergence}
\end{figure*}

% ===========================================================================
\section{Quantum circuits and depth comparison}
\label{sec:circuits}
% ===========================================================================
All quantum operations in this work are compiled to a common gate set of
single-qubit rotations and CNOTs, so the resource costs of the two routes to the
high-resolution ground state --- a traditional full-space UCCSD--VQE versus the
rodeo-terminated pipeline --- can be compared gate by gate.  Every number quoted
below is produced (and verified against exact matrix exponentials) in the
accompanying notebook.

\subsection{Circuit building blocks}
\label{sec:blocks}
The elementary primitive is the Pauli rotation $e^{-i\frac{\phi}{2}P}$ for a
Pauli string $P=\bigotimes_q\sigma_q$: each active qubit is rotated into the $Z$
eigenbasis ($H$ for $X$, $HS^\dagger$ for $Y$), a CNOT ladder accumulates the
parity onto the last active qubit, a single $R_z(\phi)$ is applied, and the
ladder and basis changes are uncomputed.  A string of weight $w$ therefore costs
$2(w-1)$ CNOTs and one parametrised rotation; the circuit reproduces the exact
exponential to $3\times10^{-16}$.

\paragraph{UCCSD preparation circuit.}
Through the Jordan--Wigner map a single excitation compiles into $2$ Pauli
rotations whose weight grows with the distance between the spin orbitals
($8$ CNOTs for adjacent levels, $16$ for the $0\!\to\!4$ single at $k=3$), and a
generic double excitation into $8$ rotations ($48$--$64$ CNOTs).  The doubles
that dominate the pairing problem are special: for a time-reversed pair moving
between levels $p$ and $q$ the operators act on the adjacent qubits
$(2p,2p{+}1)$ and $(2q,2q{+}1)$, so the internal Jordan--Wigner $Z$ strings
cancel and all $8$ Pauli strings have weight $4$ --- $48$ CNOTs per pair-hopping
term \emph{independent of the level separation}.  The full state preparation at
$N=4$, $g=1$ then costs $8$ excitation operators, $40$ rotations and $272$ CNOTs
at $k=3$, and $26$ operators, $160$ rotations and $1312$ CNOTs at $k=4$
(Table~\ref{tab:uccsd}); at $k=2$ the ansatz is empty (the sector is a single
determinant) and the preparation is gate-free.

\paragraph{Hardware-efficient ansatz.}
The HEA~\cite{Kandala2017} used for contrast consists of $L$ layers of
$R_y$ rotations on all $2k$ qubits interleaved with a CNOT entangler
(nearest-neighbour pairs plus a wrap-around bond), i.e.\ $2k(L{+}1)$ variational
angles and only $O(2kL)$ CNOTs: $12$ angles at $k=2$ ($L=2$) and $30$ at $k=3$
($L=4$).  Its depth is far smaller than UCCSD, but the circuit carries no
particle-number or seniority structure; a number-penalty term and multi-start
optimisation are needed at $k=3$, and at $k=4$ the optimiser fails to reach the
ground state (Fig.~\ref{fig:methods}).  Depth is not the only currency:
trainability of the symmetry-blind landscape becomes the bottleneck first.

\paragraph{Rodeo cycle.}
One rodeo cycle adds a Hadamard on the ancilla, the controlled evolution
$e^{-iHt_k}$, the ancilla phase $P(Et_k)$, a second Hadamard, and a mid-circuit
measurement; the ancilla is reused across cycles.  The controlled evolution is
compiled from the same Pauli-rotation primitive: in each rotation sub-circuit
only the central $R_z$ needs to be controlled (if the rotation is absent the
basis changes and CNOT ladder cancel), and a controlled $R_z$ costs $2$ extra
CNOTs.  For the pairing Hamiltonian~\eqref{eq:H} a first-order Trotter step
contains the $k(k{-}1)/2$ pair-hopping terms ($8$ weight-$4$ rotations, $48$
CNOTs each), $k$ diagonal $ZZ$ rotations ($2$ CNOTs each) and $2k$ single-qubit
phases, i.e.
\begin{equation}
  N_{\rm CNOT}^{\rm step}(k)=24\,k(k-1)+2k
  \label{eq:trottercost}
\end{equation}
CNOTs uncontrolled --- $296$ at $k=4$ --- and $\approx 4.2\times10^{2}$ CNOTs per
step once the controls are added.  A rodeo cycle with $n_T$ Trotter steps
therefore costs $\approx 4.2\times10^{2}\,n_T$ CNOTs at $k=4$.

\subsection{Depth budget: UCCSD--VQE versus the rodeo-terminated pipeline}
\label{sec:depth}
Table~\ref{tab:depth} summarises the comparison.  The traditional route prepares
the $k=4$ ground state directly with the gate-level UCCSD--VQE: every energy or
gradient evaluation inside the classical loop executes the full $1312$-CNOT,
$26$-parameter circuit, each evaluation averaged over many shots, and the
attainable accuracy saturates at the ansatz error,
$E-E_{\rm FCI}=1.4\times10^{-3}$ at $k=4$ (Table~\ref{tab:uccsd}).  The cost
also grows steeply with resolution --- $272\to1312$ CNOTs from $k=3$ to $k=4$ ---
because the operator pool grows with both $k$ and $N$.

The pipeline never runs a variational circuit at full resolution.  Its only
full-resolution circuits are (i) the parameter-free Trotter steps of the
adiabatic refinement, $296$ CNOTs each by Eq.~\eqref{eq:trottercost}, executed
once per repetition with no optimiser attached, and (ii) $M=1$--$2$ rodeo
cycles of $\approx4.2\times10^{2}\,n_T$ CNOTs each.  Even for several Trotter
steps per cycle, a rodeo cycle remains a fraction of the depth of a single
$k=4$ UCCSD preparation --- and it is executed a handful of times rather than
throughout an optimisation loop, with the post-selection acceptance held at
$p\simeq0.998$ by the refinement.  The pay-off in accuracy runs in the same
direction: two rodeo cycles reach an infidelity of $5\times10^{-6}$
(Fig.~\ref{fig:convergence}), three orders of magnitude below the variational
floor of the deeper UCCSD circuit.  The coherent depth of a single pipeline run
is dominated by the refinement stage ($n_s$ Trotter steps along the path); this
cost is fixed and parameter-free, its length is set by the physical gap through
$T\sim1/\Delta E$, and the rodeo filter is self-correcting with respect to the
residual Trotter error~\cite{Qian2024}, so both $n_s$ and $n_T$ can be kept
modest.

\begin{table}[t]
\caption{\label{tab:depth}Circuit-depth budget at $k=4$ ($N=4$, $g=1$,
eight system qubits).  Top: the traditional route, a full-space gate-level
UCCSD--VQE.  Bottom: the stages of the hybrid pipeline (one repetition; the
acceptance of the rodeo post-selection is $\approx p=0.998$).  $n_s$ and $n_T$
denote the number of Trotter steps in the refinement and in one controlled
evolution, respectively; CNOT counts follow Table~\ref{tab:uccsd} and
Eq.~\eqref{eq:trottercost}.}
\begin{ruledtabular}
\begin{tabular}{lccc}
circuit & params & CNOTs & executions\\
\colrule
UCCSD--VQE prep.\ ($k{=}4$) & 26 & 1312 & per eval.\ $\times$ shots\\
\colrule
coarse UCCSD ($k{=}2$) & 0 & 0 & 1\\
prolongation $k{=}2\!\to\!4$ & 0 & 0 & ---\\
refinement step ($k{=}4$) & 0 & 296 & $n_s$\\
rodeo cycle ($k{=}4$) & 0 & $\approx420\,n_T$ & $M=1$--$2$\\
\end{tabular}
\end{ruledtabular}
\end{table}

\subsection{Remarks on adaptive and improved VQE}
\label{sec:adapt}
The steep UCCSD gate count is not irreducible, and much recent work aims
precisely at shallower variational circuits (see Ref.~\cite{Tilly2022} for a
review).  ADAPT-VQE~\cite{Grimsley2019} grows the ansatz one operator at a
time, at each step appending the pool operator with the largest energy gradient
and re-optimising all amplitudes; it typically reaches a target accuracy with a
small fraction of the full UCCSD operator count and correspondingly shallower
circuits.  The pairing model makes the mechanism transparent: because the
Hamiltonian conserves seniority and the reference is seniority-zero, the
gradients of all singles and of all pair-breaking doubles vanish identically,
so an adaptive selection would retain only the $\frac{N}{2}(k-\frac{N}{2})$
pair-hopping doubles --- four operators at $k=4$, $N=4$, i.e.\
$4\times48=192$ CNOTs in place of the $1312$ of the full $26$-operator ansatz,
a sevenfold compression, with the discarded amplitudes expected to vanish at
the optimum just as the singles do.  Qubit-ADAPT-VQE~\cite{Tang2021} pushes
further by drawing the pool from individual Pauli strings, removing the
eightfold string overhead per fermionic double.  These developments are
complementary to the present pipeline rather than competing with it: the
pipeline is agnostic to how the coarse state is produced, so an ADAPT-style
coarse solver slots in directly, reducing the (already modest) low-resolution
cost and potentially raising the prolonged overlap; conversely, no variational
compression by itself supplies the exponential, optimiser-free convergence that
the rodeo filter provides once the overlap is large.

% ===========================================================================
\section{Discussion and outlook}
\label{sec:discussion}
% ===========================================================================
The pairing model lets us exhibit the entire EIGEN-Q pipeline with every stage
benchmarked against exact results: a gate-level UCCSD coarse solver, a gate-free
prolongation, a gapped adiabatic refinement with $T\sim1/\Delta E$, and a rodeo filter
with $1-\mathcal{F}_M\lesssim\tfrac{1-p}{p}2^{-M}$.  The two algorithms are genuinely
complementary --- refinement provides the large overlap $p\approx1$ that the rodeo
algorithm requires, and the rodeo algorithm removes the slow adiabatic tail that
refinement alone leaves behind.

Moving to hardware changes only the controlled evolution.  On a device $e^{-iHt}$ is
Trotterised over the $2k$ qubits, which introduces an effective Hamiltonian whose
eigenvalues differ from the exact ones at a controllable order in the step size; as
noted in Ref.~\cite{Qian2024}, the rodeo algorithm filters the eigenstates of this
effective Hamiltonian, and its self-correcting design (small unitary errors broaden the
energy window rather than biasing the estimate) makes it unusually robust to such
errors.  The mid-circuit measurement and deep coherent circuits required by the two
phases are supported by neutral-atom architectures such as the Magne system, making the
pipeline a natural target for a hardware demonstration.

Beyond the pairing model, the same construction applies to the targets of the EIGEN-Q
proposal --- light and medium-mass nuclei benchmarked against coupled-cluster and
in-medium similarity renormalization group calculations, molecular systems toward
chemical accuracy, and Hubbard models refined from a coarse to a fine lattice --- with
a closed-loop control layer that can adapt the refinement schedule from the rodeo
acceptance statistics read out mid-circuit.  The present results establish, on a fully
controlled example, that the hybrid of resolution refinement and the rodeo algorithm
prepares correlated eigenstates with exponential accuracy at a cost set by an overlap
that the refinement itself makes large.

% ===========================================================================
\section*{Code availability}
The complete implementation --- pairing Hamiltonian, gate-level UCCSD--VQE,
prolongation, adiabatic refinement, and the ancilla-based rodeo algorithm --- is
provided as a single Jupyter notebook that reproduces every figure in this article.

\begin{acknowledgments}
This work builds on the constant-pairing material of the FYS4480 and FYS5419 courses
at the University of Oslo and \emph{Lecture Notes in Physics} \textbf{936}.  We
acknowledge the EIGEN-Q collaboration and the QuNorth flagship programme.
\end{acknowledgments}

% ===========================================================================
\begin{thebibliography}{99}
% ===========================================================================
\bibitem{Peruzzo2014}
A.~Peruzzo, J.~McClean, P.~Shadbolt, M.-H.~Yung, X.-Q.~Zhou, P.~J.~Love,
A.~Aspuru-Guzik, J.~L.~O'Brien,
\textit{A variational eigenvalue solver on a photonic quantum processor},
Nat.\ Commun.\ \textbf{5}, 4213 (2014).

\bibitem{McClean2018}
J.~R.~McClean, S.~Boixo, V.~N.~Smelyanskiy, R.~Babbush, H.~Neven,
\textit{Barren plateaus in quantum neural network training landscapes},
Nat.\ Commun.\ \textbf{9}, 4812 (2018).

\bibitem{Albash2018}
T.~Albash, D.~A.~Lidar,
\textit{Adiabatic quantum computation},
Rev.\ Mod.\ Phys.\ \textbf{90}, 015002 (2018).

\bibitem{Abrams1999}
D.~S.~Abrams, S.~Lloyd,
\textit{Quantum algorithm providing exponential speed increase for finding
eigenvalues and eigenvectors},
Phys.\ Rev.\ Lett.\ \textbf{83}, 5162 (1999).

\bibitem{Choi2021}
K.~Choi, D.~Lee, J.~Bonitati, Z.~Qian, J.~Watkins,
\textit{Rodeo algorithm for quantum computing},
Phys.\ Rev.\ Lett.\ \textbf{127}, 040505 (2021).

\bibitem{Qian2024}
Z.~Qian, J.~Watkins, G.~Given, J.~Bonitati, K.~Choi, D.~Lee,
\textit{Demonstration of the rodeo algorithm on a quantum computer},
Eur.\ Phys.\ J.\ A \textbf{60}, 151 (2024).

\bibitem{Bogner2026}
S.~Bogner \textit{et al.},
\textit{Quantum state preparation with resolution refinement},
Phys.\ Lett.\ B \textbf{875}, 140363 (2026).

\bibitem{LeeReview2023}
D.~Lee,
\textit{Quantum techniques for eigenvalue problems},
Eur.\ Phys.\ J.\ A \textbf{59}, 275 (2023).

\bibitem{LeeUCC2019}
J.~Lee, W.~J.~Huggins, M.~Head-Gordon, K.~B.~Whaley,
\textit{Generalized unitary coupled cluster wavefunctions for quantum computation},
J.\ Chem.\ Theory Comput.\ \textbf{15}, 311 (2019).

\bibitem{Kandala2017}
A.~Kandala, A.~Mezzacapo, K.~Temme, M.~Takita, M.~Brink, J.~M.~Chow,
J.~M.~Gambetta,
\textit{Hardware-efficient variational quantum eigensolver for small molecules
and quantum magnets},
Nature \textbf{549}, 242 (2017).

\bibitem{Grimsley2019}
H.~R.~Grimsley, S.~E.~Economou, E.~Barnes, N.~J.~Mayhall,
\textit{An adaptive variational algorithm for exact molecular simulations on a
quantum computer},
Nat.\ Commun.\ \textbf{10}, 3007 (2019).

\bibitem{Tang2021}
H.~L.~Tang, V.~O.~Shkolnikov, G.~S.~Barron, H.~R.~Grimsley, N.~J.~Mayhall,
E.~Barnes, S.~E.~Economou,
\textit{Qubit-ADAPT-VQE: an adaptive algorithm for constructing
hardware-efficient ans\"atze on a quantum processor},
PRX Quantum \textbf{2}, 020310 (2021).

\bibitem{Tilly2022}
J.~Tilly, H.~Chen, S.~Cao, D.~Picozzi, K.~Setia, Y.~Li, E.~Grant, L.~Wossnig,
I.~Rungger, G.~H.~Booth, J.~Tennyson,
\textit{The variational quantum eigensolver: a review of methods and best
practices},
Phys.\ Rep.\ \textbf{986}, 1 (2022).

\bibitem{Trotter1959}
H.~F.~Trotter,
\textit{On the product of semi-groups of operators},
Proc.\ Am.\ Math.\ Soc.\ \textbf{10}, 545 (1959).

\bibitem{Suzuki1976}
M.~Suzuki,
\textit{Generalized Trotter's formula and systematic approximants of exponential
operators},
Commun.\ Math.\ Phys.\ \textbf{51}, 183 (1976).

\bibitem{Childs2021}
A.~M.~Childs, Y.~Su, M.~C.~Tran, N.~Wiebe, S.~Zhu,
\textit{Theory of Trotter error with commutator scaling},
Phys.\ Rev.\ X \textbf{11}, 011020 (2021).

\bibitem{LNP936}
M.~Hjorth-Jensen, M.~P.~Lombardo, U.~van~Kolck (eds.),
\textit{An Advanced Course in Computational Nuclear Physics},
Lecture Notes in Physics \textbf{936}, Springer (2017).

\bibitem{NielsenChuang}
M.~A.~Nielsen, I.~L.~Chuang,
\textit{Quantum Computation and Quantum Information},
Cambridge University Press (2010).
\end{thebibliography}

\end{document}
